\unitopen{15}
\begin{candos}
\item opowiedzieć o weekendzie i wczorajszym dniu w czasie przeszłym \de{Perfekt};
\item utworzyć \de{Partizip II} czasowników regularnych (\de{gemacht, gearbeitet, eingekauft, studiert});
\item zbudować zdanie z klamrą: \de{haben} na 2.\ miejscu, imiesłów na końcu;
\item zapytać o przeszłość: \de{Was hast du gemacht?}
\end{candos}

\section{Słownictwo}
\vocabtable{15}

\section{Perfekt: \de{haben} $+$ Partizip II}
W mowie i w e-mailach niemiecki opowiada o przeszłości głównie w \emph{Perfekcie}: odmieniony \de{haben} (lub \de{sein},
jednostka~16) na drugim miejscu, imiesłów na końcu. To ta sama klamra co przy czasownikach modalnych.

\begin{center}\small
\begin{tabular}{@{}lll@{}}
\toprule
\textsf{\textbf{Czasownik regularny}} & \textsf{\textbf{Partizip II}} & \textsf{\textbf{Zdanie}}\\
\midrule
\de{machen} & \de{ge}\emph{mach}\de{t} & \de{Ich habe Sport gemacht.}\\
\de{arbeiten} & \de{ge}\emph{arbeit}\de{et} & \de{Wir haben viel gearbeitet.}\\
\de{lernen} & \de{ge}\emph{lern}\de{t} & \de{Du hast Deutsch gelernt.}\\
\de{kaufen} & \de{ge}\emph{kauf}\de{t} & \de{Sie hat ein Auto gekauft.}\\
\bottomrule
\end{tabular}
\end{center}
Wzór: \de{ge-} $+$ rdzeń $+$ \de{-t} (po rdzeniu na \de{-t/-d}: \de{-et}). Brak \de{ge-} w trzech grupach:
\textbf{(1)} czasowniki na \de{-ieren} (\de{studieren $\rightarrow$ studiert}, \de{telefonieren}, \de{installieren});
\textbf{(2)} nierozdzielne przedrostki \de{be-, ge-, er-, ver-, zer-, ent-, emp-} (\de{besuchen $\rightarrow$ besucht},
\de{erklären $\rightarrow$ erklärt}); \textbf{(3)} \emph{rozdzielne} dostają \de{ge} w środku:
\de{einkaufen $\rightarrow$ ein\textbf{ge}kauft}, \de{aufräumen $\rightarrow$ auf\textbf{ge}räumt}.

\begin{zasada}[Szyk w Perfekcie]
\de{Gestern \textbf{habe} ich lange \textbf{gearbeitet}.} \de{Am Wochenende \textbf{haben} wir Freunde \textbf{besucht}.}
W pytaniu: \de{\textbf{Hast} du schon \textbf{gegessen}?} Przy \de{und} drugi czasownik pomocniczy można pominąć:
\de{Ich habe eingekauft und gekocht.}
\end{zasada}

\begin{falle}
Polski ma jeden czas przeszły („robiłem/zrobiłem”), więc kusi, by używać \emph{Präteritum}
(\de{Ich machte}). W mowie to brzmi książkowo albo niemiecko-niemiecko; wyjątki to \de{war}, \de{hatte} i czasowniki modalne
(jednostka~17). Drugi błąd: pominięcie \de{ge-}: \de{Ich habe Sport machen} (źle). I trzeci: polskie „byłem na zakupach” to
\de{Ich habe eingekauft} lub \de{Ich war einkaufen}, nie \de{Ich bin eingekauft}.
\end{falle}

\begin{szwajcaria}
W Szwajcarii mowa codzienna używa Perfektu niemal wyłącznie (dialekt nie ma Präteritum). Pisząc maila
do klienta, nadal możesz użyć Perfektu: \de{Ich habe Ihnen die Unterlagen gestern geschickt.}
W CH \de{Ich bin gestanden} i \de{Ich bin gesessen} (z \de{sein}) to regionalizm; w standardzie
\de{Ich habe gestanden / gesessen}.
\end{szwajcaria}

\begin{praca}
Status w standupie: \de{Gestern habe ich den Bug gefunden und einen Fix programmiert. Heute teste ich ihn.}
Raport: \de{Wir haben das Release geplant und die Dokumentation aktualisiert.} \de{Ich habe den Server
neu gestartet.} (\de{neu starten} jest rozdzielne: \de{neu gestartet}).
\end{praca}

\section{Dialogi}
\textbf{Dialog 1 --- w poniedziałek rano}
\begin{dialog}
\turn{A}{Wie war dein Wochenende?}{Jak minął twój weekend?}
\turn{B}{Sehr gut! Am Samstag habe ich eingekauft und gekocht, und am Sonntag habe ich Freunde besucht.}{Bardzo dobrze! W sobotę zrobiłem zakupy i gotowałem, a w niedzielę odwiedziłem przyjaciół.}
\turn{A}{Hast du auch Deutsch gelernt?}{Czy uczyłeś się też niemieckiego?}
\turn{B}{Ja, ich habe zwei Stunden gelernt und ein Video gesehen.}{Tak, uczyłem się dwie godziny i obejrzałem film.}
\turn{A}{Toll! Und was hast du gestern gemacht?}{Świetnie! A co robiłeś wczoraj?}
\end{dialog}

\section{Ćwiczenia}
\begin{cwiczenia}
\exitem{Podaj \de{Partizip II}: \de{machen}.}{\de{gemacht}}
\exitem{Podaj \de{Partizip II}: \de{arbeiten}.}{\de{gearbeitet}}
\exitem{Podaj \de{Partizip II}: \de{studieren}.}{\de{studiert}}
\exitem{Podaj \de{Partizip II}: \de{besuchen}.}{\de{besucht}}
\exitem{Podaj \de{Partizip II}: \de{einkaufen}.}{\de{eingekauft}}
\exitem{Podaj \de{Partizip II}: \de{installieren}.}{\de{installiert}}
\exitem{Utwórz zdanie w Perfekcie: \de{ich / Deutsch / lernen}.}{\de{Ich habe Deutsch gelernt.}}
\exitem{Utwórz zdanie w Perfekcie: \de{wir / das Release / planen}.}{\de{Wir haben das Release geplant.}}
\exitem{Utwórz pytanie w Perfekcie: \de{du / schon / essen}.}{\de{Hast du schon gegessen?}}
\exitem{Zacznij od \de{Gestern}: \de{ich / arbeiten / lange}.}{\de{Gestern habe ich lange gearbeitet.}}
\exitem{Popraw: \de{Ich habe Sport machen.}}{\de{Ich habe Sport gemacht.}}
\exitem{Przetłumacz: „Wczoraj testowałem kod.”}{\de{Gestern habe ich den Code getestet.}}
\exitem{Przetłumacz: „Co robiłeś w weekend?”}{\de{Was hast du am Wochenende gemacht?}}
\end{cwiczenia}

\section{Trening słownictwa (generowany)}
\drill{15}

\weekend{Opowiedz po niemiecku o minionym weekendzie w ośmiu zdaniach w Perfekcie. Zaznacz kolorem
wszystkie \de{Partizip II}. Nagraj 60 sekund i wyślij do korekty: sprawdzamy \de{ge-} i miejsce imiesłowu.}
